\documentclass[a4paper,11pt]{article}

% --- кириллица и шрифты ---
\usepackage{cmap}
\usepackage[T2A]{fontenc}
\usepackage[utf8]{inputenc}
\usepackage[russian]{babel}

\usepackage{amsmath, amsfonts, amssymb, amsthm}
\usepackage{mathrsfs}
\usepackage{geometry}
\usepackage{xcolor}

% --- пользовательские макросы, перенесённые из Beamer ---
\newcommand{\R}{\mathbb{R}}
\newcommand{\Cm}{\mathbb{C}}
\newcommand{\N}{\mathbb{N}}
\newcommand{\Z}{\mathbb{Z}}
\newcommand{\Q}{\mathbb{Q}}
\newcommand{\suml}[1]{\sum\limits_{#1=1}}
\newcommand{\Tr}{\operatorname{tr}}
\newcommand{\Ker}{\operatorname{Ker}}
\newcommand{\Img}{\operatorname{Im}}
\newcommand{\rg}{\operatorname{rg}}
\newcommand{\Con}{\mathscr{C}}
\newcommand{\Prb}{\mathbb{P}}
\newcommand{\Lp}{\mathbb{L}}
\newcommand{\E}{\mathbb{E}}
\newcommand{\Hilb}{\mathcal{H}}

\DeclareMathOperator*{\argmax}{arg\,max}

\let\eps\varepsilon
\let\Om\Omega
\let\om\omega
\let\<\langle
\let\>\rangle
\let\ph\varphi

\begin{document}

title: {Бинепрерывные полугруппы\\ и вопросы их аппроксимации}
author: {А.~В.~Уткин}
affiliation: {МИАН}
event: {Chernoff Approximations 2025}
date: {22.11.2025}
pause: 5.0


common:
line gap: {\vspace{0.8cm}}

style default: size=32 color=TEXT anim=write align=center
style sub default: size=24 color=MUTED


%--------------- СЦЕНА 1 -------

scene: {Мотивация}

line l1: {Существуют и используются полугруппы, не являющиеся сильно непрерывными.}

line l2: {Нужно работать не только в нормированных, но и в локально выпуклых пространствах.}

line l3: {\textbf{Практическая значимость:}}

line l4: {{\color{yellow}$\blacktriangleright$ Некорректные задачи Коши.}}

line l5: {{\color{yellow}$\blacktriangleright$ Сопряжённые полугруппы.}}

line l6: {{\color{yellow}$\blacktriangleright$ Динамические системы и марковские процессы.}}

line l7: {{\color{yellow}$\blacktriangleright$ Квантовая механика и квантовые каналы.}}


sub s1: {Классическая теория сильно непрерывных полугрупп слишком ограничительна для целого ряда естественных примеров. Бинепрерывные полугруппы возникают в контексте сходимости, контролируемой двумя топологиями.}


T: 0.5 s1

T: 1.5 l1

T: 1.5 l2

T: 1.0 l3

T: 1.0 l4

T: 1.0 l5

T: 1.0 l6

T: 1.0 l7

pause: 4.0


%--------------- СЦЕНА 2 -------

scene: {Бидопустимые пространства}

line l1: {Пусть $(X,\|\cdot\|)$ --- банахово пространство, а $\tau$ --- дополнительная топология на $X$.}

line l2: {\textbf{Bi.1.} $\tau$ хаусдорфова, локально выпукла и $\tau\subset\tau_{\|\cdot\|}$.}

line l3: {\textbf{Bi.2.} Ограниченные по норме $\tau$-замкнутые множества секвенциально $\tau$-полны.}

line l4: {\textbf{Bi.3.} $(X,\tau)'$ является нормирующим для $(X,\|\cdot\|)$:}

line l5: {$$\|f\|=\sup\limits_{\ph\in\Phi_X}|\langle\ph,f\rangle|.$$}

line l6: {Здесь $\Phi_X=\{\ph\in(X,\tau)' : \|\ph\|_{X^*}\le 1\}$.}


sub s1: {Бидопустимое пространство содержит в себе банахову норму и более слабую локально выпуклую топологию. Условие Bi.3 связывает сопряжённое к слабой топологии с нормой.}


T: 0.5 s1

T: 2.0 l1

T: 1.5 l2

T: 1.5 l3

T: 1.5 l4

T: 2.5 l5

T: 1.5 l6

pause: 4.0


%--------------- СЦЕНА 3 -------

scene: {Комментарии к бидопустимости}

line l1: {Топология $\tau$ порождается семейством полунорм, непрерывных по норме.}

line l2: {Через $\mathcal{P}_\tau$ обозначим все непрерывные полунормы $p$ с $p(x)\le\|x\|$.}

line l3: {Bi.2 эквивалентно тому, что равномерно ограниченные $\tau$-фундаментальные последовательности сходятся в $\tau$.}

line l4: {К Bi.3: $(X,\tau)'$ канонически вкладывается в $(X,\|\cdot\|)^*$.}

line l5: {Всякий $\tau$-непрерывный функционал непрерывен и по норме.}


sub s1: {Эти комментарии объясняют, как проверять аксиомы в конкретных пространствах. Bi.2 --- условие секвенциальной полноты, а Bi.3 управляет нормой через $\tau$-непрерывные функционалы.}


T: 0.5 s1

T: 2.0 l1

T: 1.5 l2

T: 2.0 l3

T: 3.5 l4

T: 1.5 l5

pause: 4.0


%--------------- СЦЕНА 4 -------

scene: {Пример: тривиальное бидопустимое пространство}

line l1: {\textbf{Пример 1 (тривиальный).} Пространство $(X,\|\cdot\|,\tau_{\|\cdot\|})$ бидопустимо.}

line l2: {Bi.1 очевидно: топология нормы хаусдорфова и локально выпукла.}

line l3: {Bi.2 выполнено, так как $(X,\|\cdot\|)$ --- банахово.}

line l4: {Bi.3 следует из стандартного следствия теоремы Хана--Банаха.}

line l5: {Итак, бидопустимые пространства обобщают банаховы.}


sub s1: {Простейший пример --- банахово пространство, где вторая топология совпадает с топологией нормы. Все аксиомы сводятся к классическим фактам о банаховых пространствах.}


T: 0.5 s1

T: 2.0 l1

T: 1.5 l2

T: 1.5 l3

T: 1.5 l4

T: 1.5 l5

pause: 4.0


%--------------- СЦЕНА 5 -------

scene: {Пример 2: слабая топология}

line l1: {\textbf{Пример 2.} Пусть $X$ --- банахово пространство, $\tau_w$ --- слабая топология.}

line l2: {Тогда $(X,\|\cdot\|,\tau_w)$ бидопустимо.}

line l3: {Bi.1: $\tau_w$ хаусдорфова, локально выпукла и слабее топологии нормы.}

line l4: {Bi.2: всякая ограниченная слабо фундаментальная последовательность слабо сходится.}

line l5: {Bi.3: норма восстанавливается по единичному шару:}

line l6: {$$\|x\|=\sup\limits_{\ph\in B_1^*}|\langle \ph, x \rangle|.$$}


sub s1: {Банахово пространство со слабой топологией – естественный источник нетривиальных примеров бидопостимых пространств.}


T: 0.5 s1

T: 2.0 l1

T: 1.5 l2

T: 2.5 l3

T: 1.5 l4

T: 1.5 l5

T: 2.5 l6

pause: 4.0


%--------------- СЦЕНА 6 -------

scene: {Пример 3: пространство непрерывных функций}

line l1: {\textbf{Пример 3.} Пусть $E$ --- банахово пространство, $X=C_b(E)$ с супремум-нормой.}

line l2: {Пусть $\tau_b$ --- топология сходимости на ограниченных подмножествах $E$.}

line l3: {Bi.1 проверяется легко.}

line l4: {Bi.3 следует из того, что $\delta$-функционалы $\tau_b$-непрерывны.}

line l5: {Проверка Bi.2 --- несложное упражнение.}


sub s1: {Этот пример важен для марковских полугрупп: ограниченные непрерывные функции со сходимостью на ограниченных множествах.}


T: 0.5 s1

T: 2.0 l1

T: 1.5 l2

T: 3.0 l3

T: 1.5 l4

T: 1.5 l5

pause: 4.0


%--------------- СЦЕНА 7 -------

scene: {Бисходимость и биплотность}

line l1: {Пусть $(X,\|\cdot\|,\tau)$ --- бидопустимое пространство.}

line l2: {Последовательность $(x_n)$ бисходится к $x$, если она $\|\cdot\|$-ограничена и $\tau$-сходится к $x$.}

line l3: {Подпространство $Y\subset X$ биплотно, если каждый $x\in X$\\ есть бипредел последовательности из $Y$.}

line l4: {Например, $\mathcal{S}(\R)$ биплотно в $C_b(\R)$ в топологии примера 3.}


sub s1: {Бисходимость соединяет ограниченность по норме со сходимостью в более слабой топологии. Биплотность --- соответствующее понятие аппроксимации.}


T: 0.5 s1

T: 1.5 l1

T: 2.0 l2

T: 2.0 l3

T: 2.0 l4

pause: 4.0


%--------------- СЦЕНА 8 -------

scene: {Роль условия Bi.3}

line l1: {Пусть $(x_n)$ бисходится к $x$. Тогда}

line l2: {$$\|x\|\le\sup\limits_{n\in\N}\|x_n\|.$$}

line l3: {Действительно, $\|x\|=\sup\limits_{\ph\in\Phi_X}\lim_n|\langle\ph,x_n\rangle|\le\varlimsup\limits_n\|x_n\|$.}

line l4: {Если $C=\sup\limits_n\|S_n\|<\infty$ и $S_n x$ бисходится к $Sx$ для всех $x$,}

line l5: {то $\|S\|\le C$.}


sub s1: {Условие Bi.3 даёт контроль нормы у бипределов. Это ключевой инструмент переноса оценок с приближающих операторов на предельный.}


T: 0.5 s1

T: 1.5 l1

T: 2.5 l2

T: 3.0 l3

T: 3.5 l4

T: 1.5 l5

pause: 4.5


%--------------- СЦЕНА 9 -------

scene: {Роль условия Bi.2}

line l1: {Пусть $x(\cdot):[0,T]\to X$ бинепрерывна:}

line l2: {$$\operatorname{bi}\lim\limits_{t\to t_0}x(t)=x(t_0)\quad\text{для всех }t_0\in[0,T].$$}

line l3: {Условие Bi.2 позволяет определить интеграл Римана}

line l4: {$$\int_0^T x(t)\,dt=\operatorname{bi}\lim\limits_{n\to\infty}\sum\limits_{k=1}^n x\left(\frac{kT}{n}\right)\ \dfrac{T}{n}.$$}

line l5: {Римановы суммы равномерно ограничены по норме и $\tau$-фундаментальны.}


sub s1: {Bi.2 обеспечивает секвенциальную полноту для конструкции интеграла. Именно поэтому бинепрерывные полугруппы имеют генератор и резольвенту.}


T: 0.5 s1

T: 1.5 l1

T: 2.5 l2

T: 1.5 l3

T: 3.0 l4

T: 3.5 l5

pause: 4.5


%--------------- СЦЕНА 10 -------

scene: {Контрпример: поточечная топология}

line l1: {\textbf{Вопрос.} Бидопустимо ли $X=C[-1,1]$ с супремум-нормой и топологией $\tau_p$ поточечной сходимости?}

line l2: {\textbf{Ответ:} нет.}

line l3: {Рассмотрим $f_n(y)=\arctan(ny)$. Тогда $\|f_n\|\le 1$.}

line l4: {В каждой точке $y\in[-1,1]$ последовательность $(f_n(y))$ фундаментальна, значит $\tau_p$-фундаментальна.}

line l5: {Но её поточечный предел не является непрерывным на $[-1,1]$.}


sub s1: {Этот контрпример показывает, что Bi.2 может нарушаться даже для очень естественной топологии на $C[a,b]$. Поточечная сходимость не является секвенциально полной на ограниченных множествах.}


T: 0.5 s1

T: 2.5 l1

T: 1.5 l2

T: 2.0 l3

T: 2.0 l4

T: 2.0 l5

pause: 4.5


%--------------- СЦЕНА 11 -------

scene: {Бинепрерывные полугруппы: определение}

line l1: {Полугруппа $(T_t)_{t\ge0}$ на бидопустимом пространстве называется бинепрерывной, если:}

line l2: {(i) $T_t\in\mathcal{L}(X,\|\cdot\|)$ и $\|T_t\|\le Me^{wt}$ для некоторых $M\ge1$, $w\in\R$;}

line l3: {(ii) $\operatorname{bi}\lim\limits_{t\to0+}T_t x=x$ для каждого $x\in X$;}

line l4: {(iii) если $(x_n)_n$ бисходится к $x$, то $(T_t x_n)_n$ бисходится к $T_t x$}

line l5: {и сходимость равномерна по $t$ на компактных подмножествах $\R_+$:}

line l6: {$$\sup\limits_{t\in[0,T]}p\bigl(T_t(x_n-x)\bigr)\to0\quad(n\to\infty).$$}


sub s1: {Третье условие --- ключевое требование равностепенной непрерывности. Оно заменяет сильную непрерывность равномерной бинепрерывностью на компактных промежутках времени.}


T: 0.5 s1

T: 4.0 l1

T: 2.0 l2

T: 1.5 l3

T: 2.0 l4

T: 1.5 l5

T: 2.5 l6

pause: 4.5


%--------------- СЦЕНА 12 -------

scene: {Марковские полугруппы}

line l1: {\textbf{Марковские полугруппы.} На $B_b(\R)$ с супремум-нормой оператор $T$ называется марковским, если}

line l2: {$Tf\ge0$ при $f\ge0$ и $T1=1$.}

line l3: {Тогда $0\le|Tf|\le T|f|\le\|f\|$, значит, $T$ ограничен.}

line l4: {Для однородного марковского процесса $(X_t)$ с переходными вероятностями $P_t(dy|x)$:}

line l5: {$$(T_t f)(x)=\E[f(X_t)\mid X_0=x]=\int f(y)\,P_t(dy|x).$$}


sub s1: {Марковские полугруппы --- главный источник примеров. Они ограничены на $B_b(\R)$, но часто не являются там сильно непрерывными.}


T: 0.5 s1

T: 1.5 l1

T: 1.5 l2

T: 1.5 l3

T: 1.5 l4

T: 2.5 l5

pause: 4.0


%--------------- СЦЕНА 13 -------

scene: {Разрывные марковские полугруппы}

line l1: {\textbf{Полугруппа сдвига:} $(T_t f)(y)=f(y+t)$.}

line l2: {Не сильно непрерывна на $B_b(\R)$ и $C_b(\R)$.}

line l3: {Сильно непрерывна на $\Lp_p(\R)$, $UC_b(\R)$ и $C_0(\R)$.}

line l4: {\textbf{Полугруппа теплопроводности:} $(T_t f)(y)=\int f(y+\sqrt{t}u)\,d\gamma(u)$.}

line l5: {Не сильно непрерывна на $B_b(\R)$ и $C_b(\R)$.}

line l6: {Сильно непрерывна на $\Lp_p(\R)$, $UC_b(\R)$ и $C_0(\R)$.}


sub s1: {Обе полугруппы классические. Бинепрерывная теория строится именно на тех пространствах, где они разрывны.}


T: 0.5 s1

T: 1.5 l1

T: 2.0 l2

T: 2.0 l3

T: 1.5 l4

T: 2.0 l5

T: 2.0 l6

pause: 4.5


%--------------- СЦЕНА 14 -------

scene: {Полугруппа Орнштейна--Уленбека}

line l1: {\textbf{Полугруппа Орнштейна--Уленбека:}}

line l2: {$$(T_t f)(y)=\int f(e^{-t}y+\sqrt{1-e^{-2t}}u)\,d\gamma(u).$$}

line l3: {Не сильно непрерывна на $B_b(\R)$ и $C_b(\R)$.}

line l4: {Сильно непрерывна на $\Lp_p(\R)$, $UC_b(\R)$ и $C_0(\R)$.}


sub s1: {У полугруппы Орнштейна--Уленбека та же картина: сильной непрерывности нет на пространстве всех ограниченных непрерывных функций.}


T: 0.5 s1

T: 1.5 l1

T: 2.5 l2

T: 2.0 l3

T: 2.0 l4

pause: 4.0


%--------------- СЦЕНА 15 -------

scene: {Свойство Феллера и его нарушение}

line l1: {\textbf{Свойство Феллера.} Марковская полугруппа обладает свойством Феллера,\\ если $C_0(\R)$ (или $C_b(\R)$) инвариантно.}

line l2: {Примеры выше сохраняют $C_0(\R)$ и сильно непрерывны на нём.}

line l3: {Но существуют явные контрпримеры.}

line l4: {На $\R_+$ рассмотрим $x_y(t)=\frac{y}{1+ty}$ и $(T_t f)(y)=f(x_y(t))$.}

line l5: {Эта полугруппа не сохраняет $C_0(\R_+)$.}

line l6: {$C_b(\R_+)$ инвариантно, но на нём полугруппа не является сильно непрерывной.}


sub s1: {Когда свойство Феллера нарушается, классическая теория неприменима. Естественным пространством становится $C_b$ с более слабой топологией.}


T: 0.5 s1

T: 2.0 l1

T: 1.5 l2

T: 1.0 l3

T: 2.0 l4

T: 2.0 l5

T: 2.0 l6

pause: 4.5


%--------------- СЦЕНА 16 -------

scene: {Что делать?}

line l1: {Использовать теорию бинепрерывных полугрупп.}

line l2: {\textbf{Практические приложения:}}

line l3: {{\color{yellow} $\blacktriangleright$ Полугруппы, порождённые потоками на $C_b(E)$.}}

line l4: {{\color{yellow} $\blacktriangleright$ Марковские полугруппы случайных процессов на $C_b(E)$.}}

line l5: {{\color{yellow} $\blacktriangleright$ Реализованные и сопряжённые полугруппы.}}

line l6: {{\color{yellow} $\blacktriangleright$ Квантовые динамические полугруппы на $\mathcal{B}(\Hilb)$ в слабой$^*$-топологии.}}


T: 1.5 l1

T: 1.0 l2

T: 1.5 l3

T: 1.5 l4

T: 1.5 l5

T: 1.5 l6

pause: 4.0


%--------------- СЦЕНА 17 -------

scene: {Генератор и резольвента}

line l1: {Для бинепрерывной полугруппы $(T_t)_{t\ge0}$ генератор есть $(A,D(A))$:}

line l2: {$$D(A)=\Bigl\{x\in X:\ \exists\operatorname{bi}\lim\limits_{t\to0+}\frac{T_t x-x}{t}=:Ax\Bigr\}.$$}

line l3: {При $\lambda>w$ резольвента имеет вид}

line l4: {$$R(\lambda)x=\mathcal{R}\text{-}\int_0^{+\infty}e^{-\lambda t}T_t x\,dt.$$}


sub s1: {Генератор и резольвента определяются так же, как в классической теории, только все пределы --- бипределы.}


T: 0.5 s1

T: 1.5 l1

T: 3.0 l2

T: 1.5 l3

T: 2.5 l4

pause: 4.5


%--------------- СЦЕНА 18 -------

scene: {Основы теории бинепрерывных полугрупп}

line l1: {Биплотность и бизамкнутость генератора.}

line l2: {Совпадение резольвенты полугруппы и резольвенты генератора.}

line l3: {Обобщение теоремы Хилле--Иосиды.}

line l4: {Обобщения первой и второй теорем Троттера--Като и теоремы Чернова.}

line l5: {\textbf{Теорема.} Если $X$ рефлексивно, то любая бинепрерывная полугруппа\\ на $X$ автоматически сильно непрерывна.}


sub s1: {Абстрактная теория повторяет классическую.}

sub s2: {На самом деле бинепрерывность существенно сложнее лишь в нерефлексивных пространствах.}


T: 0.5 s1

T: 1.5 l1

T: 1.5 l2

T: 1.5 l3

T: 2.0 l4

T: 3.0 l5

T: 0.5 s1

pause: 4.5


%--------------- СЦЕНА 19 -------

scene: {К теореме Чернова: равностепенная бинепрерывность}

line l1: {Семейство $(S_\alpha)_\alpha\subset\mathcal{L}(X)$ называется равностепенно бинепрерывным, если:}

line l2: {для всякой $(x_n)$, бисходящейся к $x$, $\operatorname{bi}\lim\limits_n S_\alpha x_n=x$ равномерно по $\alpha$.}

line l3: {Пусть $F_t:\R_+\to\mathcal{L}(X)$. Семейство $\{(e^{-wt}F_t)^m\}$ локально равномерно равностепенно бинепрерывно, если:}

line l4: {для всякой $(x_n)$, бисходящейся к нулю,}

line l5: {$$\tau\lim\limits_{n\to\infty}(e^{-wt}F_t)^m x_n=0,$$}

line l6: {равномерно по $t$ на компактах из $\R_+$ и по $m\in\N$.}


sub s1: {Это дополнительная равномерность, нужная, чтобы перенести итерации Чернова с биплотной области на всё пространство.}


T: 2.0 l1

T: 2.5 l2

T: 0.5 s1

T: 2.0 l3

T: 1.0 l4

T: 2.5 l5

T: 2.0 l6

pause: 4.5


%--------------- СЦЕНА 20 -------

scene: {Теорема Чернова: предположения}

line l1: {Пусть $F:[0,+\infty)\to\mathcal{L}(X)$ удовлетворяет условиям:}

line l2: {(i) $F_0=I$ и $\|F_t^m\|\le Me^{mwt}$ для некоторых $M\ge1$, $w\in\R$;}

line l3: {(ii) $\{e^{-mwt}(F_t)^m\}_{t\ge0,\ m\in\N}$ локально равномерно равностепенно бинепрерывно;}

line l4: {(iii) $\operatorname{bi}\lim\limits_{t\to0}\frac{F_t x-x}{t}=Lx$ на биплотном подмножестве $\mathcal{D}\subset X$;}

line l5: {и $(\lambda_0-L)(\mathcal{D})$ биплотно при некотором $\lambda_0>w$.}


sub s1: {Это прямой бинепрерывный аналог классической теоремы Чернова. От области требуется лишь биплотность.}


T: 0.5 s1

T: 1.5 l1

T: 2.0 l2

T: 2.5 l3

T: 2.5 l4

T: 2.0 l5

pause: 4.5


%--------------- СЦЕНА 21 -------

scene: {Теорема Чернова: заключение}

line l1: {Тогда бизамыкание $L$ есть генератор бинепрерывной полугруппы $(T_t)_{t\ge0}$.}

line l2: {Более того, для каждого $x\in X$}

line l3: {$$T_t x=\operatorname{bi}\lim\limits_{n\to\infty}(F_{t/n})^n x,$$}

line l4: {равномерно на компактных промежутках из $\R_+$.}


sub s1: {Таким образом, формула Чернова остаётся справедливой в бинепрерывной постановке. Это центральная теорема аппроксимации в приложениях.}


T: 0.5 s1

T: 2.5 l1

T: 1.0 l2

T: 2.5 l3

T: 2.0 l4

pause: 5.0


%--------------- СЦЕНА 22 -------

scene: {Вопрос 1: Ослабление условий Чернова}

line l1: {Теорема Чернова требует, чтобы $\mathcal{D}$ и $(\lambda-L)\mathcal{D}$ были биплотны в $X$.}

line l2: {Сохранится ли заключение, если они не биплотны,}

line l3: {но \textbf{если $X$ --- минимальное бизамкнутое подпространство, их содержащее}?}

line l4: {Если $\tau$ метризуема, вопрос тривиален.}

line l5: {В общем случае множество бипредельных точек подпространства может не быть бизамкнутым.}

line l6: {Минимальное бизамкнутое подпространство может быть строго больше множества бипредельных точек.}


sub s1: {Задача --- заменить биплотность более слабым относительным условием, не теряя приближения Чернова.}


T: 0.5 s1

T: 2.0 l1

T: 2.0 l2

T: 2.0 l3

T: 1.5 l4

T: 2.0 l5

T: 2.0 l6

pause: 4.5


%--------------- СЦЕНА 23 -------

scene: {Вопрос 1: Мотивация}

line l1: {Нужно, когда область определения $L=F'_0$ не биплотна в $X$.}

line l2: {Цель: найти максимальное бизамкнутое подпространство $Y\subset X$.}

line l3: {На $Y$ итерации Чернова $(F_{t/n}^n)_{n\in\N}$ должны приближать бинепрерывную полугруппу.}


sub s1: {Такая постановка естественна, когда приближения действуют только на замкнутом инвариантном подпространстве.}


T: 0.5 s1

T: 2.0 l1

T: 2.0 l2

T: 2.5 l3

pause: 4.0


%--------------- СЦЕНА 24 -------

scene: {Вопрос 1: Возможные идеи}

line l1: {1) Попытаться видоизменить теорему Троттера--Като.}

line l2: {2) Построить другую норму $|||\cdot|||$ на подпространстве $Y$, содержащем $\mathcal{D}$.}

line l3: {3) Потребовать согласованности с $(X,|||\cdot|||,\tau)$.}

line l4: {4) И сделать $\mathcal{D}$ биплотным в $(Y,|||\cdot|||,\tau)$.}


sub s1: {Более сильная вспомогательная норма могла бы восстановить биплотность на меньшем пространстве. Трудность --- в согласовании с исходной топологией.}


T: 0.5 s1

T: 1.5 l1

T: 2.0 l2

T: 1.5 l3

T: 2.0 l4

pause: 4.0


%--------------- СЦЕНА 25 -------

scene: {Вопрос 2: Аппроксимации резольвенты}

line l1: {Пусть $(F_t)_{t\ge0}$ --- функция Чернова для $(T_t)_{t\ge0}$.}

line l2: {При каких условиях верно, что для всех $x\in X$ и $\lambda>\lambda_0$}

line l3: {$$\operatorname{bi}\lim\limits_{n\to\infty}\int_0^{+\infty}e^{-\lambda t}F_{t/n}^n x\,dt=R(\lambda)x?$$}

line l4: {Здесь $R(\lambda)$ --- резольвента полугруппы $(T_t)_{t\ge0}$.}

line l5: {\textbf{Мотивация:} вычисление резольвенты, когда явной формулы для полугруппы нет.}


sub s1: {Вопрос в том, сходятся ли приближения Чернова не только к самой полугруппе, но и к её преобразованию Лапласа.}


T: 0.5 s1

T: 2.0 l1

T: 1.5 l2

T: 3.0 l3

T: 1.5 l4

T: 2.0 l5

pause: 4.5


%--------------- СЦЕНА 26 -------

scene: {Вопрос 3: Скорость сходимости}

line l1: {Обобщить теорему Галкина--Ремизова о скорости на бинепрерывные полугруппы.}

line l2: {То есть оценить $\|F_{t/n}^n x-T_t x\|$ в бинепрерывных предположениях.}

line l3: {\textbf{Мотивация:} построение быстро сходящихся итераций Чернова.}

line l4: {Основная цель --- полугруппы квантовых каналов на $\mathcal{B}(\Hilb)$.}


sub s1: {Оценки скорости нужны и для численных, и для теоретических приложений. Бинепрерывная постановка требует новых условий равномерности.}


T: 0.5 s1

T: 2.0 l1

T: 2.0 l2

T: 1.5 l3

T: 2.0 l4

pause: 4.0


%--------------- СЦЕНА 27 -------

scene: {Теорема Галкина--Ремизова: постановка}

line l1: {Пусть $F$ --- банахово пространство, $(e^{tL})_{t\ge0}$ --- $C_0$-полугруппа на $F$.}

line l2: {Зафиксируем $T>0$. Предположим:}

line l3: {(i) $\|e^{tL}\|\le M_1e^{wt}$ для всех $t\in[0,T]$, где $M_1\ge1$, $w\ge0$;}

line l4: {(ii) $S:(0,T]\to\mathcal{L}(F)$ и $\|S(t)^k\|\le M_2e^{kwt}$ для всех $k\in\N$;}

line l5: {(iii) $m\in\N_0$, $p\in\N$ и $D\subset D(L^{m+p})$ инвариантно относительно $(e^{tL})$;}

line l6: {и существуют функции $K_j:(0,T]\to[0,\infty)$, $j=0,\dots,m+p$, что:}


sub s1: {Это классические предположения теоремы Галкина--Ремизова. Область $D$ инвариантна и лежит в области определения более высокой степени генератора.}


T: 0.5 s1

T: 2.0 l1

T: 1.0 l2

T: 2.0 l3

T: 2.0 l4

T: 2.0 l5

T: 1.5 l6

pause: 4.5


%--------------- СЦЕНА 28 -------

scene: {Теорема Галкина--Ремизова: условие}

line l1: {Для всех $t\in(0,T]$ и всех $f\in D$}

line l2: {$$\Bigl\|S(t)f-\sum\limits_{k=0}^{m}\frac{t^kL^kf}{k!}\Bigr\|\le t^{m+1}\sum\limits_{j=0}^{m+p}K_j(t)\|L^jf\|.$$}

line l3: {(Это локальная оценка приближения порядка $m$)}


sub s1: {Оценка контролирует ошибку одного шага через старшие степени генератора. Это вход для оценки скорости в формуле произведения.}


T: 0.5 s1

T: 1.5 l1

T: 4.0 l2

T: 2.0 l3

pause: 5.0


%--------------- СЦЕНА 29 -------

scene: {Теорема Галкина--Ремизова: оценка отклонения}

line l1: {{\Large \textbf{Тогда:}} 1) Для всех $t>0$, $n\ge t/T$ и $f\in D$}

line l2: {$$\bigl\|S(t/n)^n f-e^{tL}f\bigr\|\le\frac{M_1M_2t^{m+1}e^{wt}}{n^m}\sum\limits_{j=0}^{m+p}C_j(t/n)\|L^jf\|.$$}

line l3: {Здесь $C_{m+1}(t)=K_{m+1}(t)e^{-wt}+\frac{M_1}{(m+1)!}$.}

line l4: {При $j\ne m+1$ имеем $C_j(t)=K_j(t)e^{-wt}$.}


sub s1: {Это явная скорость: порядок $m$ по числу итераций. Константы зависят от полугруппы и от функции Чернова.}


T: 0.5 s1

T: 1.5 l1

T: 4.0 l2

T: 2.0 l3

T: 1.5 l4

pause: 5.0


%--------------- СЦЕНА 30 -------

scene: {Теорема Галкина--Ремизова: сходимость}

line l1: {2) Пусть $D$ плотно в $F$ и $K_j(t)=o(t^{-m})$ при $t\to0^+$ для всех $j$.}

line l2: {Тогда для каждого $g\in F$ и каждого $\mathcal{T}>0$}

line l3: {$$\lim\limits_{n\to\infty,\ n\ge\mathcal{T}/T}\ \sup\limits_{t\in(0,\mathcal{T}]}\bigl\|S(t/n)^n g-e^{tL}g\bigr\|=0.$$}


sub s1: {При условиях плотности и поведения при малых $t$ формула произведения сходится равномерно на компактных промежутках времени.}


T: 0.5 s1

T: 2.5 l1

T: 1.5 l2

T: 3.5 l3

pause: 5.0


%--------------- СЦЕНА 31 -------

scene: {Литература}

line l1: {[1] F.~Kühnemund. Bi-continuous\ semigroups on spaces with two topologies. Ph.D. thesis, Tübingen, 2001.}

line l2: {[2] A.A.~Albanese, E.~Mangino. Trotter--Kato theorems for bi-continuous semigroups. J. Math. Anal. Appl. 289 (2004), 477--492.}

line l3: {[3] I.D.~Remizov. Chernoff approximations and resolvents of linear operators. arXiv:2301.06765 (2023).}

line l4: {[4] O.E.~Galkin, I.D.~Remizov. Rate of convergence in the Chernoff product formula. Isr. J. Math. 265 (2025), 929--943.}


sub s1: {В этих работах изложены основы теории бинепрерывных полугрупп и результаты об аппроксимациях Чернова, обсуждавшиеся в докладе.}


T: 0.5 s1

T: 2.0 l1

T: 2.5 l2

T: 2.5 l3

T: 2.5 l4

pause: 5.0


%--------------- СЦЕНА 32 -------

scene: {}

line l1: {\Huge Спасибо за внимание!}

line l2: {\Large Вопросы?}


T: 0.0 l1

T: 2.0 l2

pause: 4.0

\end{document}